% ECONOMICS_PROBLEM: {"id": "L11-310", "title": "Reliable markup measurement", "jel_code": "L11", "source_id": "OP-0391"}
% !TeX program = xelatex
\documentclass[11pt,a4paper]{article}
\usepackage[margin=23mm]{geometry}
\usepackage{fontspec}
\setmainfont{Latin Modern Roman}
\usepackage{amsmath,amssymb,mathtools}
\usepackage{xcolor,enumitem,booktabs,longtable,array,xurl,hyperref}
\definecolor{muted}{HTML}{53636C}
\hypersetup{colorlinks=true,urlcolor=blue,linkcolor=blue}
\setlength{\parindent}{0pt}
\setlength{\parskip}{5pt}
\newcommand{\R}{\mathbb R}
\newcommand{\E}{\mathbb E}
\newcommand{\N}{\mathbb N}
\newcommand{\ind}{\mathbf 1}
\newcommand{\Var}{\operatorname{Var}}
\newcommand{\Cov}{\operatorname{Cov}}
\newcommand{\supp}{\operatorname{supp}}
\DeclareMathOperator*{\argmin}{arg\,min}
\DeclareMathOperator*{\argmax}{arg\,max}
\newcommand{\entrymeta}[3]{{\sffamily\small\color{muted}\textbf{\texttt{#1}}\quad|\quad #2\par #3\par}}
\newcommand{\Problemref}[1]{\texttt{#1}}
\newcommand{\JELref}[2]{\texttt{#2} (JEL \texttt{#1})}
\begin{document}
\section*{Reliable markup measurement}
\textbf{Problem ID:} \texttt{L11-310}\quad \textbf{JEL:} \texttt{L11}\par
% BEGIN ECONOMICS STATEMENT
\label{problem:OP-0391}
\label{audited:ECON-032}
\label{audited:ECON-035}
{\sffamily\small\color{muted}Original subject group: Industrial organization and competition\par}
\label{definitions:problem:30007673}
\textbf{Audit classification:} Published research agenda. The source motivates validation, not this specific partial-identification theorem. Expenditures alone do not determine marginal costs; a lower cost bound is needed for finite markup bounds.
\subsection*{Definitions and precise research target}
Consider multiproduct manufacturing firms selling differentiated products in United States markets during a fixed ten-year observation window. For product \(j\) in period \(t\), observe price \(p_{jt}>0\), quantity \(q_{jt}\), materials expenditure, labor expenditure and total firm expenditure on intangible inputs. Let \(C_t(q_t,k_t;\theta)\) be differentiable total variable cost, where \(q_t\) is the firm's vector of product quantities, \(k_t\) its intangible-input stock and \(\theta\) unknown cost parameters. Product marginal cost is \(c_{jt}=\partial C_t/\partial q_{jt}>0\), and the target markup is \(\mu_{jt}=p_{jt}/c_{jt}\). Shared expenditures cannot be assigned to products without restrictions. Let \(\mathcal C(P,\mathcal R)\) denote all cost functions consistent with observed-data distribution \(P\) and a stated restriction set \(\mathcal R\), such as independently measured output elasticities and input-price variation. For prespecified nonnegative weights \(w_{jt}\) summing to one, bound the aggregate target \(\sum_{j,t}w_{jt}\mu_{jt}\) over \(\mathcal C(P,\mathcal R)\). Identify additional product-cost observations that shrink those bounds. Require that any point estimator be checked against independent engineering or accounting measurements unavailable during estimation. The problem concerns reliable measurement within this population, with uncertainty about cost allocation reported explicitly rather than absorbed into an assumed pricing model.

\textbf{Definitions, assumptions and mathematical qualifications.} For each observed firm-period output vector \(q\), a candidate cost model includes input prices, cost shocks, the observed expenditure measurement equation, and an independently stated stock law for intangible capital. \(\mathcal C(P,\mathcal R)\) is the nonempty collection of such complete models reproducing \(P\); a bare cost function does not generate a data distribution. Assume \(0<c_{\min}\leq\partial_j C\leq c_{\max}<\infty\) on observed support. Fixed weights \(w_{jt}\) are determined before counterfactual changes, and \(A(C)=\sum_{jt}w_{jt}p_{jt}/\partial_jC\). Sharp endpoints are \(\inf_{C\in\mathcal C}A(C)\) and \(\sup_{C\in\mathcal C}A(C)\); endpoint attainment is not assumed. A validation audit measures an independently defined marginal production cost, rather than average accounting cost or an allocation of fixed costs.
\subsection*{Variants and assumption changes}
\begin{enumerate}
\item \textbf{Editorial, benchmark.} For a single product with an externally observed differentiable variable-cost curve, marginal cost and markup are known on observed support; this is a solved measurement benchmark.
\item \textbf{Editorial, joint inputs.} Permit nonseparable \(C(q,k)\), unknown allocation of common input expenditures and bounded audit error \(|c^{\rm audit}_{jt}-c_{jt}|\leq e_{jt}\). Compare sharp aggregate bounds with and without this audit.
\end{enumerate}
\subsection*{References and source audit}
\textbf{Checked source.} 24 August 2024 author draft, Section 6, printed p. 36, third research direction. Full primary text retrieved. \url{https://web.stanford.edu/~ayurukog/trendsincompetition.pdf}.
\begin{enumerate}\raggedright
\item\hypertarget{definitions:source:30007673:1}{} Carl Shapiro; Ali Yurukoglu. 2024. \emph{Trends in Competition in the United States: What Does the Evidence Show?}. 24 August 2024 author draft, Section 6, printed p. 36, third research direction.
\url{https://web.stanford.edu/~ayurukog/trendsincompetition.pdf}.
DOI: \url{https://doi.org/10.3386/w32762}.
\end{enumerate}

\subsection*{Integrated refinement: ECON-032}
{\sffamily\small\color{muted}Markup measurement with input wedges and production uncertainty\par}
This refinement adopts the additional source conventions
(page~\pageref{audited:conventions}), subject to its local restrictions.
\textbf{Flexible-input proxies with endogenous input wedges.}
Let physical output be $q=F(V,K,L;\omega)$,
$\theta^V=\partial\log F/\partial\log V$ and $s^V=p_VV/(pq)$.
The production proxy $\widetilde\mu=\theta^V/s^V$ need not equal
the markup $\mu=p/MC$. With effective marginal input cost $(1+\xi)p_V$,
an interior cost-minimization condition instead gives
$\widetilde\mu=\mu(1+\xi)$. Characterize the compatible markup set over a declared
class of input wedges and measurement/production models reproducing the observed
law. Require positive prices, inputs and marginal costs, differentiability,
appropriate moments and an explicit weighting population. Distinguish physical
from price-contaminated revenue elasticities, and firm weights from sales weights.
Unrestricted wedges may leave the set uninformative; specify which independent
cost or input-price data shrink it.
\subsection*{Supplied source audit for ECON-032}
Local [S1], [S2], etc. in this audit and the references immediately below
belong to ECON-032, independently of the original entry's reference numbering.
[S1] and [S2], especially Sections II.B and IV of the checked working version, support measurement concerns. The identified set and example wedge are editorial.
\subsection*{References for integrated refinement ECON-032}
{\footnotesize\raggedright
\par\noindent\textbf{[S1]} Chad Syverson (2019). \emph{Macroeconomics and Market Power: Context, Implications, and Open Questions}. Journal of Economic Perspectives 33(3), 23–43. DOI: 10.1257/jep.33.3.23.\par \textit{Verified locator / source role:} Markup-measurement discussion and unresolved measurement questions.\par \url{https://www.aeaweb.org/articles?id=10.1257/jep.33.3.23}\par\smallskip
\par\noindent\textbf{[S2]} Chad Syverson (2019). \emph{Macroeconomics and Market Power: Facts, Potential Explanations, and Open Questions}. Brookings Economic Studies, January 2019; author research essay.\par \textit{Verified locator / source role:} Sections II.B and IV, printed pp. 9–12 and 20–21: markup measurement.\par \url{https://www.brookings.edu/wp-content/uploads/2019/01/ES_20190116_Syverson-Macro-Micro-Market-Power.pdf}\par\smallskip
}

\subsection*{Integrated refinement: ECON-035}
{\sffamily\small\color{muted}Intangible investment versus rising measured markups\par}
This refinement adopts the additional source conventions
(page~\pageref{audited:conventions}), subject to its local restrictions.
\textbf{Intangible-capital accounting versus true markup changes.}
Specify $K^I_{j,t+1}=(1-\delta_I)K^I_{jt}+I^I_{jt}$ and the measurement equation
for expensed or capitalized intangible expenditure. For fixed population weights,
identify the terms of
\[
\Delta\mathbb E[\widetilde\mu]
=\Delta\mathbb E[\mu]+\Delta\mathbb E[\widetilde\mu-\mu].
\]
The equality is an accounting identity; identifying its components is the
research target. Specify the estimator, physical output, depreciation and
production restrictions. Fixed overhead, variable intangible inputs and durable
capital can imply different marginal costs. Capitalizing an expense does not
automatically correct a production-based markup, and an intangible user cost
enters marginal cost only under the stated technology. Separate matched-firm
trends from changes in firm composition and accounting margins from
price-to-marginal-cost markups.
\subsection*{Supplied source audit for ECON-035}
Local [S1], [S2], etc. in this audit and the references immediately below
belong to ECON-035, independently of the original entry's reference numbering.
[S1] and [S2] discuss intangibles in broader investment and concentration patterns. They do not state this precise markup-bias decomposition as an open conjecture. It is an editorial extension of that background.
\subsection*{References for integrated refinement ECON-035}
{\footnotesize\raggedright
\par\noindent\textbf{[S1]} Chad Syverson (2019). \emph{Macroeconomics and Market Power: Context, Implications, and Open Questions}. Journal of Economic Perspectives 33(3), 23–43. DOI: 10.1257/jep.33.3.23.\par \textit{Verified locator / source role:} Potential explanations involving intangible capital, fixed costs, and measurement.\par \url{https://www.aeaweb.org/articles?id=10.1257/jep.33.3.23}\par\smallskip
\par\noindent\textbf{[S2]} Chad Syverson (2019). \emph{Macroeconomics and Market Power: Facts, Potential Explanations, and Open Questions}. Brookings Economic Studies, January 2019; author research essay.\par \textit{Verified locator / source role:} Section III, printed pp. 16–17: intangibles and efficiency; Section IV.\par \url{https://www.brookings.edu/wp-content/uploads/2019/01/ES_20190116_Syverson-Macro-Micro-Market-Power.pdf}\par\smallskip
}

\subsection*{Common conventions: Source mathematical conventions}

These conventions apply to each entry unless explicitly replaced.

\textbf{Models and identification.} A primitive model \(m\) specifies all probability laws, feasible choices, information, equilibrium and measurement rules used by its target. The admissible class \(\mathcal M\) is fixed independently of outcomes. If \(P_O(m)\) is its observable probability law and \(\tau(m)\) the target, its identified set at observable law \(P\) is
\[
 \mathcal I_\tau(P)=\{\tau(m):m\in\mathcal M,\ P_O(m)=P\}.
\]
Point identification means that this set is a singleton. An empty set signals incompatibility of data and assumptions. Coordinate bounds need not describe the joint set. Bounds are sharp only if every retained target is attainable or arbitrarily approachable by compatible models. Finite bounds require explicit restrictions on unobserved quantities; unrestricted confounding, hidden wealth, extrapolation or arbitrary equilibrium selection can yield uninformative sets.

\textbf{Causal interventions.} A potential outcome \(Y_i(p)\) is the outcome under a complete policy regime \(p\), including timing, financing, information and market spillovers. The target population and units are fixed before treatment. With interference, use the complete assignment vector or a justified exposure mapping; individual no-interference notation alone cannot identify an economy-wide intervention. A difference of mean potential outcomes does not require the cross-world joint distribution. Conditioning on post-treatment migration, survival, enrollment or employment changes the estimand and needs separate justification.

\textbf{Statistical statements.} An estimator, confidence set, forecast or test specifies a sampling law, model class and loss/coverage criterion. Uniform \((1-\alpha)\)-coverage means \(\inf_{m\in\mathcal M}\Pr_m\{\tau(m)\in C_n\}\geq1-\alpha\), or an explicitly asymptotic version. The whole parameter space trivially has coverage, so informativeness requires a stated width, risk or power objective. A forecasting improvement is relative to a named benchmark, information set and evaluation population.

\textbf{Equilibrium and optimization.} An equilibrium comprises feasible plans, optimality, beliefs where needed, budget constraints and all market/resource clearing. The primitives or a stated benchmark define each correspondence \(\mathcal E(p;m)\). Nonemptiness is assumed or proved, never inferred just from compact policy sets. A selected-equilibrium target uses a declared selection rule; a robust target quantifies over all permitted equilibria. Use suprema or infima unless attainment is established. An \(\varepsilon\)-optimal policy is feasible and within \(\varepsilon\) of the finite optimum value. Report an empty feasible set separately.

\textbf{Welfare and accounting.} Normative utility, interpersonal normalization, social weights, numeraire, discounting and the use of public receipts are primitives. Behavioral utility need not coincide with the welfare criterion. Household welfare includes ownership income and rents once; adding profits again to compensating variation generally double-counts them. A monetary equivalent is defined by an explicit transfer and price convention, with monotonicity and range conditions. Average effects, mechanism contributions, transfers and welfare losses are different targets. Infinite sums require convergence and dynamic feasible plans require terminal or no-Ponzi conditions.

\textbf{Source versions.} A working-paper date, revision date and journal publication year are distinguished. A DOI for a journal version is not automatically the DOI of an earlier manuscript. Locators refer to the stated version and printed pages where specified. A metadata-only check does not verify a claim about a full-text conclusion. Every entry's source subsection narrows the provenance claim accordingly.

\subsection*{Common conventions: Additional-source status and mathematical conventions}

\label{audited:conventions}

Of the 100 supplied ECON entries, 65 distinct problems appear here; 32 close
matches refine earlier entries, and three duplicate questions map to existing
entries. Appendix~\ref{app:audited-crosswalk} lists every destination.
The attachment's P/R/S/U distinctions and source audits are retained. Its
precise mathematical formalizations are editorial unless the individual audit
says otherwise. Candidate subsidy targets ECON-004--006 retain their U flags;
the resolved approval-core question ECON-007 updates OP-0202 above.
The following supplied conventions govern both this part and the attributed
ECON refinements elsewhere in the catalogue, subject to local restrictions.

\textbf{Parameterized questions.} A problem family lists inputs that must be fixed before an instance is evaluated: population, horizon, policies, information, data law, model class and welfare weights. These are required choices, not quantities the citations are assumed to specify. Undefined applications should not be treated as identified estimands.

\textbf{Indices and integrability.} Sets of agents, firms and regions are finite unless a population measure is explicitly used. \(i,j\) index units, \(r\) regions, \(t\) dates and \(h\) horizons. \(\mathbb E\) and \(\Pr\) refer to the declared population and law; \(\mathbf1\{A\}\) is an indicator. Every displayed expectation and discounted sum needs existence in the stated domain. Infinite horizons use a discount in \((0,1)\) and a convergence condition; finite horizons may use discount one.

\textbf{Optimization and existence.} A displayed \(\max\), \(\min\), or \(\arg\max\) is conditional on attainment. For finite feasible menus this is immediate; general spaces need continuity and compactness/coercivity or another existence argument. Without attainment, the target is the supremum/infimum and arbitrarily near-optimal feasible choices. \(\inf\varnothing=+\infty\). Empty feasibility and an unidentified target are different issues.

\textbf{Potential outcomes.} \(Y_i(z)\) is the outcome under a specified intervention, not the observed conditional mean among people choosing \(z\). In binary comparisons \(\Delta Y_i=Y_i(1)-Y_i(0)\). Specify assignment, treatment versions, compliance, information and observation horizon. Interference is allowed under a full policy vector; compressed exposure notation needs a justified mapping. No interference, consistency and overlap are substantive assumptions, not automatic consequences of notation.

\textbf{Populations and observation.} Fix baseline eligibility and population weights. Conditioning on post-policy employment, survival, relocation or program uptake can change the population being compared. Death is not ordinary loss to follow-up; outcomes truncated by death need a composite convention, joint survival target, or explicitly defined survivor stratum. Fertility-changing interventions need a population functional rather than an unexplained fixed descendant cohort. Measurement and missingness models must be supplied.

\textbf{Identification.} A structural model \(m\in\mathcal M\) implies observable law \(P_m\) and target \(\theta(m)\). Identification means
\[
 P_{m_1}=P_{m_2}\ \Longrightarrow\ \theta(m_1)=\theta(m_2)
 \quad\text{for all }m_1,m_2\in\mathcal M .
\]
Without identification, the target is
\[
 \Theta(P;\mathcal M)=\{\theta(m):m\in\mathcal M,\ P_m=P\}.
\]
Writing \(\theta(P)\) before establishing identification can hide incompatible latent models. Confidence sets additionally need a sampling law, confidence level, asymptotic sequence and uniformity class. Point identification, coverage and informative precision are separate properties.

\textbf{Mediation.} Total effects of randomized assignment are distinct from cross-world natural indirect effects. Belief/effort-channel effects require a meaningful mediator intervention and additional measurement, exclusion or mediation assumptions. Randomization of the initial policy alone does not supply those assumptions. Report bounds or interventional alternatives when the stronger target is unsupported.

\textbf{Equilibrium.} A counterfactual \(W(z)\), \(p(z)\), or \(\mu_t^z\) requires individual optimization, feasibility, government/resource budgets, market clearing and state-transition laws. State equilibrium existence and selection, or report ranges over compatible equilibria. Initial-state integration includes subsequent endogenous transitions; current-state integration must use the policy-dependent distribution. Reduced-form invariance after policy is not assumed.

\textbf{Welfare accounts.} Scores, utility, health, dollars and ecological quantities require explicit conversion before addition. Fees, taxes, subsidies, wages and extortion payments are transfers between parties; distributional weights can make them welfare relevant, but their face value is not automatically a resource loss. State ownership, geographic boundaries, foreign-income weights and fiscal finance. Do not add profits to household welfare already including dividends, or count land capitalization and the underlying benefit twice. A benefit-cost expression must distinguish gross benefits from net welfare and subtract every real cost exactly once.

\textbf{Derivatives and risk.} Log derivatives require positive arguments; local responses require admissible perturbations and specified equilibrium branches. Differentiation under expectations needs regularity/domination, and boundaries may allow only one-sided derivatives. Risk uses a specified law; ambiguity uses a declared nonempty law set on a common history space. Adaptive policies must be nonanticipative. A robust criterion is an editorial choice unless attributed explicitly.

\textbf{Scope overlaps.} Allocation, collective choice, contracts, household bargaining and coordinated policy overlap economics with optimization and game theory. They are retained because they appeared in the original catalogue; no claim of a strictly game-theory-free selection is made.
% END ECONOMICS STATEMENT
\end{document}
